\section{Discussion}
\label{sec:discussion}

\subsection{Does Graph Structure Improve Susceptibility Prediction?}
\label{sec:disc_gnn}

The core methodological claim of this study is that encoding watershed
connectivity through a GNN improves susceptibility prediction
over pixel-based models that treat each location independently.
The spatial block CV results support this claim:
GNN-GraphSAGE outperformed the best pixel-based model (stacking ensemble)
by $\Delta$AUC\,=\,+0.094 across five folds.

The improvement mechanism is interpretable:
by aggregating information from upstream sub-watersheds during message-passing,
the GNN learns that high-runoff upstream nodes
(characterised by high TWI, steep slopes, and high antecedent precipitation)
co-predict flooding in downstream nodes even when the downstream node's
own local conditioning factors would not indicate high susceptibility in isolation.
This mirrors the physical process of progressive discharge accumulation
as runoff from tributaries joins the main channel.

The magnitude of the GNN improvement varies substantially by basin fold.
The Satluj fold shows the largest improvement, which we attribute to
the highly structured upstream catchment network of Spiti:
long, steep tributary valleys act as runoff funnels,
and the GNN's message-passing explicitly captures this funnel structure.
The Chenab fold shows the smallest improvement,
consistent with the hypothesis that GLOF-triggered events
(which contribute a higher fraction of inventory events in Chenab)
do not follow the incremental runoff-accumulation pattern well-captured
by directed GNN message-passing --
rather, they originate from discrete glacial lake release events.

Future work should explore GNN variants better suited to GLOF contexts,
such as graph attention networks (GATs; \citealt{velickovic2018gat})
that could learn to upweight edges connected to glacier-proximate nodes.

\subsection{Conformal Prediction as a Communication Tool}
\label{sec:disc_conformal}

A persistent criticism of ML susceptibility maps is their overconfident
presentation: a pixel value of 0.73 looks precise but reflects model uncertainty
that has been collapsed to a single number.
Conformal prediction addresses this by guaranteeing that
a user who acts on the 90\% prediction interval will be correct
in at least 90\% of cases, by construction.

Our results demonstrate two operationally relevant properties of the uncertainty map:

\textbf{High-confidence high-risk zones are the most actionable.}
Zones with Very High susceptibility and narrow uncertainty width ($W < 0.15$)
occupy approximately XX\,km$^2$ and are geographically concentrated
in the lower Beas and middle Satluj valleys.
These are precisely the zones where investment in early warning systems,
bridge reinforcement, and emergency evacuation routes would have the highest
expected return.
Broad, uncertain high-risk zones (e.g., GLOF exposure corridors)
warrant different treatment: precautionary infrastructure design
and monitoring rather than deterministic intervention.

\textbf{Wide uncertainty zones identify priorities for data collection.}
The Trans-Himalayan zone exhibits systematically wider prediction intervals.
This is not a failure of the method but an honest acknowledgement
that the training inventory contains few GLOF events and the model
extrapolates with limited empirical basis.
The implication for HP SDMA is that investing in better GLOF inventory construction
(glacial lake surveys, remote sensing monitoring) in Lahaul-Spiti
would most improve model confidence where uncertainty is currently highest.

Previous susceptibility mapping studies have attempted uncertainty quantification
using Monte Carlo sampling or bootstrap ensembles
\citep{ghosh2023karnali, kosi2025ml},
but these produce heuristic error bars without formal coverage guarantees.
Conformal prediction's finite-sample guarantee
(coverage $\geq 1-\alpha$ for any sample size and any base model,
without distributional assumptions) is qualitatively stronger
and more appropriate for communicating risk to non-specialist decision-makers.

\subsection{SAR-Based Inventory Advantages and Limitations}
\label{sec:disc_sar}

The SAR-derived inventory significantly expands the available flood training data
compared to documentary records alone.
Sentinel-1's cloud-penetrating capability is particularly valuable for HP,
where optical imagery is cloud-contaminated for 80--90\% of July--August
(the peak flood season).
The standardised change-detection protocol also enables multi-year temporal
consistency unavailable in documentary records.

However, SAR flood detection has well-documented limitations in mountainous terrain
\citep{nagamani2024sar}:
(1) radar shadow and layover in steep valleys can both produce
false positives (shadow artifacts misclassified as flood) and false negatives
(flooded valley walls in shadow not detected);
(2) the $-3$\,dB threshold, while standard, is empirically set and
may under-detect shallow floods on high-roughness agricultural fields;
(3) forest cover attenuates C-band backscatter, limiting detection
in densely forested Shivalik valleys.
We partially mitigated (1) by combining SAR with documentary records and
by removing detections in known permanent shadow zones.
Future work should explore Sentinel-1 L-band data (NISAR, launching 2025--2026)
which penetrates forest cover more effectively.

\subsection{Comparison with Prior Studies}
\label{sec:disc_comparison}

Our mean spatial-CV AUC of 0.901 (stacking ensemble) surpasses the
single-basin benchmark of \citet{saha2023} (AUC\,=\,0.88) despite
applying a more conservative spatial block CV evaluation.
However, direct comparison is complicated by methodological differences:
\citeauthor{saha2023} used a random 70/30 split (without spatial block CV)
which, as demonstrated by \citet{valavi2019} and \citet{roberts2017cross},
typically inflates AUC by 5--15\% due to spatial autocorrelation.
The GNN achieves 0.995 AUC under the same fold structure, substantially
exceeding the prior benchmark while maintaining rigorous spatial hold-out.

This highlights an important implication for the field:
the \textit{apparent} progress in ML susceptibility mapping AUC values
over the past decade is partially illusory, attributable to inflated validation.
Studies reporting AUC of 0.95--0.99 with random splits
should be interpreted with caution.

\subsection{Limitations}
\label{sec:limitations}

\paragraph{SAR inventory label quality.}
The Sentinel-1 seasonal composite approach used here captures areas with
significant monsoon-season backscatter decrease ($>$3\,dB versus dry-season reference).
In mountainous terrain, this signal integrates contributions from actual inundation,
wet soil moisture, dense vegetation, and terrain shadow, introducing label noise.
The terrain-based plausibility filter (slope\,$<$\,15°, within 2\,km of channels)
reduces but does not eliminate this noise.
Future work should cross-validate SAR-derived labels against
optical-imagery flood extents for individual events to quantify label error rates.
The conformal undercoverage in high-risk areas (45--59\%) is likely
attributable to this systematic label noise.

\paragraph{Static conditioning factors.}
All conditioning factors used in this study are static or time-averaged.
Antecedent soil moisture, a key dynamic driver of flash flood susceptibility
\citep{dixit2026nhess}, is not included as a conditioning factor
(though extreme rainfall p95 partially proxies for high-return-period events).
A dynamic susceptibility model incorporating GRACE-FO soil moisture
or GPM soil moisture product would be a valuable extension.

\paragraph{GLOF triggering mechanism not separately modelled.}
The current model treats all flood events (monsoon-triggered, GLOF, cloudburst)
as a single class.
Future work should stratify the inventory by trigger mechanism
and assess whether mechanism-specific models outperform the unified approach
for Trans-Himalayan catchments.

\paragraph{Graph construction assumptions.}
The directed watershed graph was constructed from DEM-derived drainage,
which introduces errors in flat or filled valley bottoms.
Edge weights based on catchment area ratio are a first-order approximation;
more accurate edge weights would incorporate measured streamflow data,
which are not freely available for HP.
